\chapter{Unit 4: Environmental Pollution and Waste Management}

\section{Section A: Multiple Choice Practice Questions (60 Curated Problems)}

\noindent\textbf{Q1. What is environmental pollution?} \hfill [Tag: Definition, types and causes of environmental pollution | Easy]
\begin{enumerate}[label=(\alph*),topsep=2pt,itemsep=1pt,partopsep=0pt,parsep=0pt]
    \item {The harmful contamination of the environment}
    \item {The seasonal change in weather conditions}
    \item {The natural renewal of environmental resources}
    \item {The planned conservation of wildlife habitats}
\end{enumerate}
\vspace{2mm}

\noindent\textbf{Q2. Which activity is a major human cause of environmental pollution?} \hfill [Tag: Definition, types and causes of environmental pollution | Easy]
\begin{enumerate}[label=(\alph*),topsep=2pt,itemsep=1pt,partopsep=0pt,parsep=0pt]
    \item {Planting native trees}
    \item {Burning fossil fuels}
    \item {Harvesting rainwater}
    \item {Protecting wetlands}
\end{enumerate}
\vspace{2mm}

\noindent\textbf{Q3. Which pollutant is commonly released by motor vehicles?} \hfill [Tag: Effects and controls of air pollution | Easy]
\begin{enumerate}[label=(\alph*),topsep=2pt,itemsep=1pt,partopsep=0pt,parsep=0pt]
    \item {Sodium chloride}
    \item {Carbon monoxide}
    \item {Dissolved oxygen}
    \item {Calcium carbonate}
\end{enumerate}
\vspace{2mm}

\noindent\textbf{Q4. Which measure helps control air pollution from factories?} \hfill [Tag: Effects and controls of air pollution | Easy]
\begin{enumerate}[label=(\alph*),topsep=2pt,itemsep=1pt,partopsep=0pt,parsep=0pt]
    \item {Installing filters in chimneys}
    \item {Increasing the use of coal}
    \item {Removing trees near factories}
    \item {Burning waste in open areas}
\end{enumerate}
\vspace{2mm}

\noindent\textbf{Q5. Which is a common cause of water pollution?} \hfill [Tag: Effects and controls of water pollution | Easy]
\begin{enumerate}[label=(\alph*),topsep=2pt,itemsep=1pt,partopsep=0pt,parsep=0pt]
    \item {Collection of rainwater}
    \item {Planting trees near rivers}
    \item {Discharge of untreated sewage}
    \item {Use of solar energy}
\end{enumerate}
\vspace{2mm}

\noindent\textbf{Q6. What is an effective method for controlling water pollution?} \hfill [Tag: Effects and controls of water pollution | Easy]
\begin{enumerate}[label=(\alph*),topsep=2pt,itemsep=1pt,partopsep=0pt,parsep=0pt]
    \item {Dumping waste farther from shore}
    \item {Treating wastewater before discharge}
    \item {Mixing sewage with river water}
    \item {Removing plants from riverbanks}
\end{enumerate}
\vspace{2mm}

\noindent\textbf{Q7. Excessive use of which substance can pollute agricultural soil?} \hfill [Tag: Effects and controls of soil pollution | Easy]
\begin{enumerate}[label=(\alph*),topsep=2pt,itemsep=1pt,partopsep=0pt,parsep=0pt]
    \item {Crop residues}
    \item {Chemical pesticides}
    \item {Green manure}
    \item {Organic compost}
\end{enumerate}
\vspace{2mm}

\noindent\textbf{Q8. Which practice helps reduce soil pollution?} \hfill [Tag: Effects and controls of soil pollution | Easy]
\begin{enumerate}[label=(\alph*),topsep=2pt,itemsep=1pt,partopsep=0pt,parsep=0pt]
    \item {Using organic manure}
    \item {Dumping electronic waste}
    \item {Applying excess fertilizer}
    \item {Burning plastic waste}
\end{enumerate}
\vspace{2mm}

\noindent\textbf{Q9. Long-term exposure to loud noise can directly cause which health problem?} \hfill [Tag: Effects and controls of noise pollution | Easy]
\begin{enumerate}[label=(\alph*),topsep=2pt,itemsep=1pt,partopsep=0pt,parsep=0pt]
    \item {Bone fracture}
    \item {Hearing loss}
    \item {Skin infection}
    \item {Tooth decay}
\end{enumerate}
\vspace{2mm}

\noindent\textbf{Q10. Which measure can reduce traffic noise in residential areas?} \hfill [Tag: Effects and controls of noise pollution | Easy]
\begin{enumerate}[label=(\alph*),topsep=2pt,itemsep=1pt,partopsep=0pt,parsep=0pt]
    \item {Widening loudspeaker use}
    \item {Increasing horn use}
    \item {Removing roadside trees}
    \item {Installing noise barriers}
\end{enumerate}
\vspace{2mm}

\noindent\textbf{Q11. Which is a possible effect of excessive exposure to ionizing radiation?} \hfill [Tag: Effects and controls of radiation pollution | Easy]
\begin{enumerate}[label=(\alph*),topsep=2pt,itemsep=1pt,partopsep=0pt,parsep=0pt]
    \item {Faster water purification}
    \item {Damage to body cells}
    \item {Improved hearing ability}
    \item {Increased soil fertility}
\end{enumerate}
\vspace{2mm}

\noindent\textbf{Q12. Which substance is commonly classified as an emerging pollutant?} \hfill [Tag: Emerging pollutants | Easy]
\begin{enumerate}[label=(\alph*),topsep=2pt,itemsep=1pt,partopsep=0pt,parsep=0pt]
    \item {Pharmaceutical residue}
    \item {Clean river sand}
    \item {Atmospheric oxygen}
    \item {Natural water vapor}
\end{enumerate}
\vspace{2mm}

\noindent\textbf{Q13. What is a major environmental effect of fireworks?} \hfill [Tag: Fireworks and their ill-effects | Easy]
\begin{enumerate}[label=(\alph*),topsep=2pt,itemsep=1pt,partopsep=0pt,parsep=0pt]
    \item {Increased forest cover}
    \item {Improved soil quality}
    \item {Air and noise pollution}
    \item {Groundwater recharge}
\end{enumerate}
\vspace{2mm}

\noindent\textbf{Q14. Which gas is the main greenhouse gas released by burning fossil fuels?} \hfill [Tag: Global warming | Easy]
\begin{enumerate}[label=(\alph*),topsep=2pt,itemsep=1pt,partopsep=0pt,parsep=0pt]
    \item {Hydrogen}
    \item {Helium}
    \item {Nitrogen}
    \item {Carbon dioxide}
\end{enumerate}
\vspace{2mm}

\noindent\textbf{Q15. Which is a likely effect of climate change?} \hfill [Tag: Climate change | Easy]
\begin{enumerate}[label=(\alph*),topsep=2pt,itemsep=1pt,partopsep=0pt,parsep=0pt]
    \item {Complete removal of storms}
    \item {Equal rainfall in all regions}
    \item {Permanent weather stability}
    \item {More frequent heat waves}
\end{enumerate}
\vspace{2mm}

\noindent\textbf{Q16. Which group of chemicals is mainly associated with ozone layer depletion?} \hfill [Tag: Ozone layer depletion | Easy]
\begin{enumerate}[label=(\alph*),topsep=2pt,itemsep=1pt,partopsep=0pt,parsep=0pt]
    \item {Amino acids}
    \item {Chlorofluorocarbons}
    \item {Natural sugars}
    \item {Carbohydrates}
\end{enumerate}
\vspace{2mm}

\noindent\textbf{Q17. Which gases are mainly responsible for the formation of acid rain?} \hfill [Tag: Acid rain and impacts on human communities and agriculture | Easy]
\begin{enumerate}[label=(\alph*),topsep=2pt,itemsep=1pt,partopsep=0pt,parsep=0pt]
    \item {Methane and argon}
    \item {Oxygen and hydrogen}
    \item {Sulfur dioxide and nitrogen oxides}
    \item {Helium and neon}
\end{enumerate}
\vspace{2mm}

\noindent\textbf{Q18. The Bhopal gas tragedy was caused by the accidental release of which toxic gas?} \hfill [Tag: Case studies on pollution | Easy]
\begin{enumerate}[label=(\alph*),topsep=2pt,itemsep=1pt,partopsep=0pt,parsep=0pt]
    \item {Carbon dioxide}
    \item {Nitrogen gas}
    \item {Methyl isocyanate}
    \item {Water vapor}
\end{enumerate}
\vspace{2mm}

\noindent\textbf{Q19. What do the three Rs of solid waste management represent?} \hfill [Tag: Solid waste management | Easy]
\begin{enumerate}[label=(\alph*),topsep=2pt,itemsep=1pt,partopsep=0pt,parsep=0pt]
    \item {Remove, replace, and refill}
    \item {Recover, repair, and return}
    \item {Record, report, and review}
    \item {Reduce, reuse, and recycle}
\end{enumerate}
\vspace{2mm}

\noindent\textbf{Q20. Which measure is suitable for controlling industrial wastewater pollution?} \hfill [Tag: Control measures of urban and industrial waste | Easy]
\begin{enumerate}[label=(\alph*),topsep=2pt,itemsep=1pt,partopsep=0pt,parsep=0pt]
    \item {Diluting waste with rainwater}
    \item {Using an effluent treatment plant}
    \item {Releasing waste into nearby drains}
    \item {Storing waste beside a river}
\end{enumerate}
\vspace{2mm}

\noindent\textbf{Q21. A residential area develops breathing problems, contaminated wells, and declining soil fertility after a chemical factory begins operating nearby. Which conclusion best describes this situation?} \hfill [Tag: Definition, types and causes of environmental pollution | Medium]
\begin{enumerate}[label=(\alph*),topsep=2pt,itemsep=1pt,partopsep=0pt,parsep=0pt]
    \item {The factory is causing pollution across multiple environmental media}
    \item {The factory is causing only occupational pollution inside its premises}
    \item {The changes are caused only by increased residential water use}
    \item {The changes are examples of natural environmental variation}
\end{enumerate}
\vspace{2mm}

\noindent\textbf{Q22. A city records high concentrations of fine particulate matter near a coal-fired power plant. Which measure would most directly reduce particulate emissions from the plant's chimney?} \hfill [Tag: Effects and controls of air pollution | Medium]
\begin{enumerate}[label=(\alph*),topsep=2pt,itemsep=1pt,partopsep=0pt,parsep=0pt]
    \item {Increasing the chimney height}
    \item {Building sound barriers nearby}
    \item {Installing an electrostatic precipitator}
    \item {Adding chlorine to cooling water}
\end{enumerate}
\vspace{2mm}

\noindent\textbf{Q23. During winter, a city experiences severe smog when cool air becomes trapped below a layer of warm air. Why does pollution increase under this condition?} \hfill [Tag: Effects and controls of air pollution | Medium]
\begin{enumerate}[label=(\alph*),topsep=2pt,itemsep=1pt,partopsep=0pt,parsep=0pt]
    \item {Warm air forces pollutants into the groundwater}
    \item {Vertical mixing of polluted air is restricted}
    \item {Wind speed increases near the ground}
    \item {Sunlight immediately destroys all gaseous pollutants}
\end{enumerate}
\vspace{2mm}

\noindent\textbf{Q24. Untreated sewage enters a lake, followed by an algal bloom and a large fish die-off. Which sequence best explains the fish mortality?} \hfill [Tag: Effects and controls of water pollution | Medium]
\begin{enumerate}[label=(\alph*),topsep=2pt,itemsep=1pt,partopsep=0pt,parsep=0pt]
    \item {Algae remove sediments, and fish lose spawning grounds}
    \item {Nutrients increase salinity, and fish lose body heat}
    \item {Nutrients increase algae, and decomposition lowers dissolved oxygen}
    \item {Sewage raises water clarity, and fish become exposed}
\end{enumerate}
\vspace{2mm}

\noindent\textbf{Q25. Wastewater from a food-processing unit has a high biochemical oxygen demand (BOD). Which treatment stage is most suitable for reducing its BOD before discharge?} \hfill [Tag: Effects and controls of water pollution | Medium]
\begin{enumerate}[label=(\alph*),topsep=2pt,itemsep=1pt,partopsep=0pt,parsep=0pt]
    \item {Mechanical grit removal}
    \item {Simple chlorination alone}
    \item {Coarse screening alone}
    \item {Aerobic biological treatment}
\end{enumerate}
\vspace{2mm}

\noindent\textbf{Q26. A farmer repeatedly applies a persistent pesticide. Years later, residues are detected in the soil and in birds feeding near the farm. Which practice would best reduce further contamination while maintaining pest control?} \hfill [Tag: Effects and controls of soil pollution | Medium]
\begin{enumerate}[label=(\alph*),topsep=2pt,itemsep=1pt,partopsep=0pt,parsep=0pt]
    \item {Applying the pesticide before every rainfall}
    \item {Burning crop residues after each harvest}
    \item {Increasing the pesticide dose each season}
    \item {Adopting integrated pest management}
\end{enumerate}
\vspace{2mm}

\noindent\textbf{Q27. Soil near a mining site contains toxic concentrations of lead. Which remediation method uses living plants to remove or stabilize the metal?} \hfill [Tag: Effects and controls of soil pollution | Medium]
\begin{enumerate}[label=(\alph*),topsep=2pt,itemsep=1pt,partopsep=0pt,parsep=0pt]
    \item {Electrostatic precipitation}
    \item {Thermal inversion}
    \item {Ultraviolet disinfection}
    \item {Phytoremediation}
\end{enumerate}
\vspace{2mm}

\noindent\textbf{Q28. A hospital is located beside a busy highway. Which intervention would most effectively reduce the noise reaching patients without changing the hospital's location?} \hfill [Tag: Effects and controls of noise pollution | Medium]
\begin{enumerate}[label=(\alph*),topsep=2pt,itemsep=1pt,partopsep=0pt,parsep=0pt]
    \item {Replacing asphalt with reflective concrete}
    \item {Constructing a dense barrier between the road and hospital}
    \item {Increasing traffic speed during visiting hours}
    \item {Installing taller streetlights along the highway}
\end{enumerate}
\vspace{2mm}

\noindent\textbf{Q29. Workers handling a sealed radioactive source want to reduce their radiation dose. Which combination follows the standard principles of radiation protection?} \hfill [Tag: Effects and controls of radiation pollution | Medium]
\begin{enumerate}[label=(\alph*),topsep=2pt,itemsep=1pt,partopsep=0pt,parsep=0pt]
    \item {Reduce exposure time, increase distance, and use shielding}
    \item {Reduce exposure time, reduce distance, and remove shielding}
    \item {Increase exposure time, reduce distance, and improve ventilation}
    \item {Increase exposure time, increase distance, and add humidity}
\end{enumerate}
\vspace{2mm}

\noindent\textbf{Q30. Trace amounts of antibiotics are detected downstream from hospitals and livestock farms. What is the most significant long-term ecological concern?} \hfill [Tag: Emerging pollutants | Medium]
\begin{enumerate}[label=(\alph*),topsep=2pt,itemsep=1pt,partopsep=0pt,parsep=0pt]
    \item {Selection and spread of antibiotic-resistant microorganisms}
    \item {Rapid neutralization of every dissolved nutrient}
    \item {Immediate conversion of freshwater into saltwater}
    \item {Permanent elimination of all aquatic microorganisms}
\end{enumerate}
\vspace{2mm}

\noindent\textbf{Q31. A city wants to reduce the short-term pollution peak during a major festival. Which policy would address both air and noise pollution from fireworks most directly?} \hfill [Tag: Fireworks and their ill-effects | Medium]
\begin{enumerate}[label=(\alph*),topsep=2pt,itemsep=1pt,partopsep=0pt,parsep=0pt]
    \item {Increasing street lighting during celebrations}
    \item {Washing roads only after the festival ends}
    \item {Closing schools for one additional day}
    \item {Restricting fireworks by time, location, and permitted type}
\end{enumerate}
\vspace{2mm}

\noindent\textbf{Q32. A power company replaces a coal-fired plant with solar generation while supplying the same amount of electricity. Why can this help limit global warming?} \hfill [Tag: Global warming | Medium]
\begin{enumerate}[label=(\alph*),topsep=2pt,itemsep=1pt,partopsep=0pt,parsep=0pt]
    \item {It increases atmospheric reflection by producing more smoke}
    \item {It removes all existing greenhouse gases from the atmosphere}
    \item {It converts atmospheric methane directly into oxygen}
    \item {It reduces carbon dioxide released during electricity generation}
\end{enumerate}
\vspace{2mm}

\noindent\textbf{Q33. A coastal farming community faces more frequent saltwater intrusion into wells and fields. Which response is primarily an adaptation to climate change?} \hfill [Tag: Climate change | Medium]
\begin{enumerate}[label=(\alph*),topsep=2pt,itemsep=1pt,partopsep=0pt,parsep=0pt]
    \item {Introducing a carbon tax on heavy industries}
    \item {Restoring forests to absorb atmospheric carbon dioxide}
    \item {Replacing coal plants with wind farms nationwide}
    \item {Developing salt-tolerant crops and protecting freshwater storage}
\end{enumerate}
\vspace{2mm}

\noindent\textbf{Q34. A refrigeration company replaces chlorofluorocarbons (CFCs) with refrigerants that do not release chlorine in the stratosphere. Which environmental problem is most directly addressed?} \hfill [Tag: Ozone layer depletion | Medium]
\begin{enumerate}[label=(\alph*),topsep=2pt,itemsep=1pt,partopsep=0pt,parsep=0pt]
    \item {Formation of ground-level dust}
    \item {Accumulation of lead in soil}
    \item {Destruction of stratospheric ozone}
    \item {Nutrient enrichment of lakes}
\end{enumerate}
\vspace{2mm}

\noindent\textbf{Q35. A farming region downwind of several power plants experiences increasingly acidic rainfall. Which combined response would address both the source and an immediate soil impact?} \hfill [Tag: Acid rain and impacts on human communities and agriculture | Medium]
\begin{enumerate}[label=(\alph*),topsep=2pt,itemsep=1pt,partopsep=0pt,parsep=0pt]
    \item {Install noise barriers and irrigate fields with saline water}
    \item {Raise chimney heights and apply additional nitrogen fertilizer}
    \item {Install flue-gas desulfurization and apply lime to acidic soil}
    \item {Use taller buildings and remove organic matter from soil}
\end{enumerate}
\vspace{2mm}

\noindent\textbf{Q36. A factory releases methylmercury into a bay. Residents who regularly eat contaminated fish develop neurological disorders. Which process best explains the high exposure through fish consumption?} \hfill [Tag: Case studies on pollution | Medium]
\begin{enumerate}[label=(\alph*),topsep=2pt,itemsep=1pt,partopsep=0pt,parsep=0pt]
    \item {Dilution at each successive trophic level}
    \item {Sedimentation of all mercury before uptake}
    \item {Neutralization by dissolved oxygen}
    \item {Biomagnification through the aquatic food chain}
\end{enumerate}
\vspace{2mm}

\noindent\textbf{Q37. Following a toxic gas leak from an industrial facility near a densely populated settlement, which measure would most directly reduce casualties during a similar future incident?} \hfill [Tag: Case studies on pollution | Medium]
\begin{enumerate}[label=(\alph*),topsep=2pt,itemsep=1pt,partopsep=0pt,parsep=0pt]
    \item {A larger parking area beside chemical storage tanks}
    \item {An off-site emergency plan with alarms and evacuation routes}
    \item {A public park located next to the factory boundary}
    \item {A taller administrative building within the facility}
\end{enumerate}
\vspace{2mm}

\noindent\textbf{Q38. A municipality collects food scraps, recyclable packaging, and hazardous household waste in one truck. Which change is most important for improving material recovery and treatment?} \hfill [Tag: Solid waste management | Medium]
\begin{enumerate}[label=(\alph*),topsep=2pt,itemsep=1pt,partopsep=0pt,parsep=0pt]
    \item {Compress all mixed waste before collection}
    \item {Burn mixed waste without emission controls}
    \item {Require segregation of waste at the source}
    \item {Transport all waste directly to open dumps}
\end{enumerate}
\vspace{2mm}

\noindent\textbf{Q39. A town has a large amount of separately collected food and garden waste. Which management option can produce both useful energy and a nutrient-rich residue?} \hfill [Tag: Solid waste management | Medium]
\begin{enumerate}[label=(\alph*),topsep=2pt,itemsep=1pt,partopsep=0pt,parsep=0pt]
    \item {Anaerobic digestion}
    \item {Deep-well injection}
    \item {Open dumping}
    \item {Uncontrolled burning}
\end{enumerate}
\vspace{2mm}

\noindent\textbf{Q40. An electroplating factory produces wastewater containing chromium and other heavy metals. Which approach best supports effective pollution control and resource recovery?} \hfill [Tag: Control measures of urban and industrial waste | Medium]
\begin{enumerate}[label=(\alph*),topsep=2pt,itemsep=1pt,partopsep=0pt,parsep=0pt]
    \item {Release the stream during periods of high river flow}
    \item {Segregate the metal-rich stream and treat it for metal recovery}
    \item {Mix the stream with domestic sewage to reduce its concentration}
    \item {Store the untreated stream in an unlined evaporation pond}
\end{enumerate}
\vspace{2mm}

\noindent\textbf{Q41. A power station releases chemically clean cooling water into a river. Downstream, dissolved oxygen decreases and temperature-sensitive fish disappear, although no toxic chemical exceeds its regulatory limit. Which interpretation is most accurate?} \hfill [Tag: Definition, types and causes of environmental pollution | Hard]
\begin{enumerate}[label=(\alph*),topsep=2pt,itemsep=1pt,partopsep=0pt,parsep=0pt]
    \item {The discharge is habitat loss because pollution requires a persistent synthetic substance}
    \item {The discharge is thermal pollution because it causes an adverse environmental change}
    \item {The discharge is not pollution because its chemical composition remains within limits}
    \item {The discharge is water pollution only if pathogens increase above background levels}
\end{enumerate}
\vspace{2mm}

\noindent\textbf{Q42. A city in a valley experiences severe particulate pollution on clear winter mornings. Emissions remain nearly constant, winds are weak, and temperature initially increases with altitude. Which process-control combination best addresses the episode?} \hfill [Tag: Effects and controls of air pollution | Hard]
\begin{enumerate}[label=(\alph*),topsep=2pt,itemsep=1pt,partopsep=0pt,parsep=0pt]
    \item {Temperature inversion traps pollutants; temporary emission reduction limits accumulation}
    \item {Convective instability traps pollutants; taller chimneys eliminate regional exposure}
    \item {Adiabatic cooling traps pollutants; water spraying removes all gaseous contaminants}
    \item {Photochemical inversion destroys pollutants; nighttime traffic restriction increases mixing}
\end{enumerate}
\vspace{2mm}

\noindent\textbf{Q43. An industrial boiler emits coarse fly ash, fine particulate matter, sulfur dioxide, and nitrogen oxides. Which control train most appropriately targets these pollutants in sequence?} \hfill [Tag: Effects and controls of air pollution | Hard]
\begin{enumerate}[label=(\alph*),topsep=2pt,itemsep=1pt,partopsep=0pt,parsep=0pt]
    \item {Electrostatic precipitator, biofilter, cyclone, and chlorination reactor}
    \item {Cyclone, electrostatic precipitator, wet scrubber, and selective catalytic reduction}
    \item {Wet scrubber, cyclone, activated carbon, and gravitational settling chamber}
    \item {Selective catalytic reduction, cyclone, sedimentation tank, and fabric aerator}
\end{enumerate}
\vspace{2mm}

\noindent\textbf{Q44. A treatment plant removes 90\% of suspended solids but little dissolved organic matter from sewage. The effluent is visually clear, yet fish kills occur downstream. Which measurement would most directly reveal the remaining oxygen-depleting pollution?} \hfill [Tag: Effects and controls of water pollution | Hard]
\begin{enumerate}[label=(\alph*),topsep=2pt,itemsep=1pt,partopsep=0pt,parsep=0pt]
    \item {Electrical conductivity measured before removal of suspended solids}
    \item {Turbidity measured immediately after gravitational sedimentation}
    \item {Biochemical oxygen demand measured over a standardized incubation period}
    \item {Total hardness measured through calcium carbonate equivalence}
\end{enumerate}
\vspace{2mm}

\noindent\textbf{Q45. A phosphorus-limited lake receives sewage rich in both nitrate and phosphate. Managers can fund only one immediate intervention. Which action is most likely to reduce harmful algal blooms at their source?} \hfill [Tag: Effects and controls of water pollution | Hard]
\begin{enumerate}[label=(\alph*),topsep=2pt,itemsep=1pt,partopsep=0pt,parsep=0pt]
    \item {Install tertiary phosphorus removal at the sewage treatment plant}
    \item {Apply algicide repeatedly while leaving sewage inputs unchanged}
    \item {Stock herbivorous fish without reducing nutrient concentrations}
    \item {Aerate deep water without changing external nutrient loading}
\end{enumerate}
\vspace{2mm}

\noindent\textbf{Q46. Two soils contain the same total mass of cadmium. Soil X is acidic with low organic matter, while Soil Y is neutral with high clay and organic matter. Why may crops grown in Soil X accumulate more cadmium?} \hfill [Tag: Effects and controls of soil pollution | Hard]
\begin{enumerate}[label=(\alph*),topsep=2pt,itemsep=1pt,partopsep=0pt,parsep=0pt]
    \item {Low pH converts cadmium completely into a harmless elemental form}
    \item {Low pH and weak sorption increase the bioavailable cadmium fraction}
    \item {High organic matter in Soil Y accelerates cadmium volatilization}
    \item {High clay content in Soil Y increases cadmium uptake through roots}
\end{enumerate}
\vspace{2mm}

\noindent\textbf{Q47. At a workplace, two independent machines each produce 90 dB and a third produces 87 dB at the same location. What is the approximate combined sound level?} \hfill [Tag: Effects and controls of noise pollution | Hard]
\begin{enumerate}[label=(\alph*),topsep=2pt,itemsep=1pt,partopsep=0pt,parsep=0pt]
    \item {91 dB}
    \item {267 dB}
    \item {94 dB}
    \item {97 dB}
\end{enumerate}
\vspace{2mm}

\noindent\textbf{Q48. A radioactive source has an activity of 640 Bq and a half-life of 8 days. After 24 days, shielding that blocks 75\% of the emitted radiation is installed. What activity-equivalent exposure remains outside the shielding?} \hfill [Tag: Effects and controls of radiation pollution | Hard]
\begin{enumerate}[label=(\alph*),topsep=2pt,itemsep=1pt,partopsep=0pt,parsep=0pt]
    \item {80 Bq}
    \item {160 Bq}
    \item {20 Bq}
    \item {40 Bq}
\end{enumerate}
\vspace{2mm}

\noindent\textbf{Q49. A municipal plant detects persistent per- and polyfluoroalkyl substances in treated wastewater. Conventional activated sludge produces little removal. Which strategy is technically most defensible?} \hfill [Tag: Emerging pollutants | Hard]
\begin{enumerate}[label=(\alph*),topsep=2pt,itemsep=1pt,partopsep=0pt,parsep=0pt]
    \item {Add nutrients to stimulate complete mineralization in the secondary clarifier}
    \item {Use primary sedimentation because dissolved PFAS readily settle without sorbents}
    \item {Increase aeration until the carbon-fluorine bonds undergo rapid biodegradation}
    \item {Use adsorption or membrane separation and manage the concentrated residual separately}
\end{enumerate}
\vspace{2mm}

\noindent\textbf{Q50. Air monitoring during a fireworks festival shows a sharp increase in fine particles, barium, strontium, and perchlorate residues. Which conclusion best connects these observations to likely impacts?} \hfill [Tag: Fireworks and their ill-effects | Hard]
\begin{enumerate}[label=(\alph*),topsep=2pt,itemsep=1pt,partopsep=0pt,parsep=0pt]
    \item {Perchlorate rapidly becomes harmless oxygen before reaching soil or surface water}
    \item {Combustion products can worsen respiratory exposure and contaminate deposited dust or runoff}
    \item {Metal salts remain entirely inside the firework casing and cannot enter environmental media}
    \item {Fine particles affect visibility but are too large to penetrate the human respiratory system}
\end{enumerate}
\vspace{2mm}

\noindent\textbf{Q51. Suppose doubling atmospheric carbon dioxide produces a radiative forcing of 3.7 W m − 2 , while aerosols contribute − 1.2 W m − 2 . If the equilibrium sensitivity parameter is 0.8 ∘ C / ( W m − 2 ) , what warming is estimated from the net forcing?} \hfill [Tag: Global warming | Hard]
\begin{enumerate}[label=(\alph*),topsep=2pt,itemsep=1pt,partopsep=0pt,parsep=0pt]
    \item {3.9 ∘}
    \item {2.0 ∘}
    \item {3.0 ∘}
    \item {1.0 ∘}
\end{enumerate}
\vspace{2mm}

\noindent\textbf{Q52. A coastal agricultural district faces rising mean sea level, more frequent storm surges, and increasing soil salinity. Which policy package combines adaptation with mitigation?} \hfill [Tag: Climate change | Hard]
\begin{enumerate}[label=(\alph*),topsep=2pt,itemsep=1pt,partopsep=0pt,parsep=0pt]
    \item {Relocate farms inland, remove coastal wetlands, and subsidize fossil-fuel irrigation}
    \item {Construct seawalls, expand diesel pumping, and increase cultivation in exposed floodplains}
    \item {Deepen drainage channels, increase groundwater extraction, and discontinue emissions monitoring}
    \item {Restore mangroves, adopt salt-tolerant crops, and reduce regional greenhouse gas emissions}
\end{enumerate}
\vspace{2mm}

\noindent\textbf{Q53. Why can chlorofluorocarbons remain comparatively harmless in the lower atmosphere yet cause substantial ozone loss in the stratosphere?} \hfill [Tag: Ozone layer depletion | Hard]
\begin{enumerate}[label=(\alph*),topsep=2pt,itemsep=1pt,partopsep=0pt,parsep=0pt]
    \item {They persist long enough to reach the stratosphere, where ultraviolet radiation releases catalytic chlorine}
    \item {They dissolve in rainwater, which transports reactive chlorine upward into polar clouds}
    \item {They dissociate in the troposphere, where chlorine permanently binds all atmospheric oxygen}
    \item {They absorb infrared radiation at the surface and directly convert ozone into carbon dioxide}
\end{enumerate}
\vspace{2mm}

\noindent\textbf{Q54. A farming region downwind of coal-fired power plants experiences soil acidification, declining crop yield, and elevated soluble aluminum. Which response most directly addresses both the cause and the damaged soil?} \hfill [Tag: Acid rain and impacts on human communities and agriculture | Hard]
\begin{enumerate}[label=(\alph*),topsep=2pt,itemsep=1pt,partopsep=0pt,parsep=0pt]
    \item {Increase ammonia fertilizer use and irrigate fields with untreated acidic runoff}
    \item {Raise chimney heights and apply additional aluminum-containing soil amendments}
    \item {Replace crop rotation with monoculture and remove all organic residues}
    \item {Reduce sulfur and nitrogen oxide emissions and apply lime based on soil testing}
\end{enumerate}
\vspace{2mm}

\noindent\textbf{Q55. In the Minamata pollution case, why did severe neurological disease occur in communities consuming fish even when mercury concentrations in surrounding water were not always exceptionally high?} \hfill [Tag: Case studies on pollution | Hard]
\begin{enumerate}[label=(\alph*),topsep=2pt,itemsep=1pt,partopsep=0pt,parsep=0pt]
    \item {Methylmercury bioaccumulated in organisms and biomagnified through the aquatic food web}
    \item {Inorganic mercury increased water hardness and indirectly caused short-term oxygen depletion}
    \item {Mercury remained confined to sediment and therefore reached people mainly through inhalation}
    \item {Elemental mercury was created naturally in fish and diluted at higher trophic levels}
\end{enumerate}
\vspace{2mm}

\noindent\textbf{Q56. A city compares two waste policies: mixed-waste incineration with energy recovery, or source separation followed by recycling, composting, and treatment of residuals. Under the waste hierarchy, when is the second policy generally preferable?} \hfill [Tag: Solid waste management | Hard]
\begin{enumerate}[label=(\alph*),topsep=2pt,itemsep=1pt,partopsep=0pt,parsep=0pt]
    \item {When organic waste is mixed with hazardous waste before centralized mechanical sorting}
    \item {When recycling consumes more resources than virgin production across its full life cycle}
    \item {When recovered materials and compost displace virgin resources with acceptable contamination}
    \item {When all separated materials must be transported farther and ultimately landfilled untreated}
\end{enumerate}
\vspace{2mm}

\noindent\textbf{Q57. An industrial estate sends variable acidic, alkaline, metal-bearing, and organic waste streams into one municipal sewer. Which management change provides the strongest preventive control?} \hfill [Tag: Control measures of urban and industrial waste | Hard]
\begin{enumerate}[label=(\alph*),topsep=2pt,itemsep=1pt,partopsep=0pt,parsep=0pt]
    \item {Increase sewer flow velocity so contaminants reach the treatment plant faster}
    \item {Require source segregation, process-specific pretreatment, and monitored discharge limits}
    \item {Dilute peak discharges with drinking water until visual pollution disappears}
    \item {Combine all streams to average their concentrations before municipal treatment}
\end{enumerate}
\vspace{2mm}

\noindent\textbf{Q58. A metal-finishing plant wants to reduce hazardous sludge rather than merely dispose of it. Which redesign best applies cleaner-production principles?} \hfill [Tag: Control measures of urban and industrial waste | Hard]
\begin{enumerate}[label=(\alph*),topsep=2pt,itemsep=1pt,partopsep=0pt,parsep=0pt]
    \item {Dilute spent baths, discharge them intermittently, and expand off-site storage}
    \item {Add precipitation chemicals earlier, maximize sludge formation, and landfill the solids}
    \item {Increase rinse-water flow, mix all wastes, and dewater the resulting sludge}
    \item {Use counter-current rinsing, recover bath chemicals, and substitute less toxic inputs}
\end{enumerate}
\vspace{2mm}

\noindent\textbf{Q59. A community has elevated indoor radon originating from soil gas beneath buildings. Which intervention most directly reduces long-term inhalation exposure?} \hfill [Tag: Effects and controls of radiation pollution | Hard]
\begin{enumerate}[label=(\alph*),topsep=2pt,itemsep=1pt,partopsep=0pt,parsep=0pt]
    \item {Increase indoor humidity and coat external walls with acoustic insulation}
    \item {Apply activated carbon to drinking water and close roof openings}
    \item {Add dense window curtains and reduce mechanical ventilation rates}
    \item {Install sub-slab depressurization and seal major foundation entry routes}
\end{enumerate}
\vspace{2mm}

\noindent\textbf{Q60. Which lesson from the 1984 Bhopal gas disaster most directly addresses low-probability but catastrophic industrial pollution risks?} \hfill [Tag: Case studies on pollution | Hard]
\begin{enumerate}[label=(\alph*),topsep=2pt,itemsep=1pt,partopsep=0pt,parsep=0pt]
    \item {Locate emergency controls only inside the plant because off-site exposure is unlikely}
    \item {Maintain layered safety systems, minimize hazardous inventories, and plan community warnings}
    \item {Rely primarily on odor detection and assume nearby residents can evacuate without planning}
    \item {Store larger reactive inventories so production can continue during maintenance shutdowns}
\end{enumerate}
\vspace{2mm}


\section{Section B: Comprehensive Subjective Problem with Detailed Analytical Rubric}

\begin{subjectiveq}{4}{1}
\textbf{Industrial Effluent River Assimilation and Streeter-Phelps Oxygen Sag:}\\A municipal-industrial combined wastewater stream is discharged continuously into a perennial freshwater river. The hydrodynamic and water quality parameters at the outfall are:\begin{itemize}[leftmargin=1.5em]  \item Wastewater Flow: $Q_w = 1.2\text{ m}^3/\text{s}$, with ultimate $\text{BOD}_u = 180\text{ mg/L}$, and $\text{DO}_w = 1.0\text{ mg/L}$.  \item River Flow: $Q_r = 10.8\text{ m}^3/\text{s}$, with upstream $\text{BOD}_u = 3.0\text{ mg/L}$, and $\text{DO}_r = 8.5\text{ mg/L}$.  \item Saturated Dissolved Oxygen at stream temperature ($22^\circ\text{C}$): $\text{DO}_{\text{sat}} = 8.80\text{ mg/L}$.  \item Deoxygenation rate constant: $k_d = 0.25\text{ day}^{-1}$ (base $e$).  \item Atmospheric Reaeration rate constant: $k_r = 0.60\text{ day}^{-1}$ (base $e$).  \item Average stream flow velocity: $u = 0.75\text{ m/s} = 64.8\text{ km/day}$.\end{itemize}(a) Calculate the initial mixed dissolved oxygen ($\text{DO}_0$), initial oxygen deficit ($D_0$), and initial ultimate BOD ($L_0$) immediately downstream of complete mixing.\\(b) Compute the critical travel time $t_c$ (in days) to reach the minimum dissolved oxygen concentration using the classical Streeter-Phelps formulation.\\(c) Compute the critical oxygen deficit $D_c$ and the minimum dissolved oxygen concentration $\text{DO}_{\text{min}}$. Verify whether $\text{DO}_{\text{min}}$ violates the Central Pollution Control Board (CPCB) minimum standard of $4.0\text{ mg/L}$ for the propagation of aquatic life and fisheries.\\(d) Determine the downstream location $x_c$ (in kilometers) from the discharge point where the minimum dissolved oxygen occurs.
\end{subjectiveq}

\begin{solbox}
\textbf{Part (a): Mass Balance at Instantaneous Mixing:}\\Total mixed stream discharge $Q_m = Q_w + Q_r = 1.2 + 10.8 = 12.0\text{ m}^3/\text{s}$.\begin{align*}\text{DO}_0 &= \frac{Q_w \text{DO}_w + Q_r \text{DO}_r}{Q_m} = \frac{(1.2 \times 1.0) + (10.8 \times 8.5)}{12.0} = \frac{1.2 + 91.8}{12.0} = \frac{93.0}{12.0} = 7.75\text{ mg/L} \\L_0 &= \frac{Q_w L_w + Q_r L_r}{Q_m} = \frac{(1.2 \times 180) + (10.8 \times 3.0)}{12.0} = \frac{216 + 32.4}{12.0} = \frac{248.4}{12.0} = 20.70\text{ mg/L} \\D_0 &= \text{DO}_{\text{sat}} - \text{DO}_0 = 8.80 - 7.75 = 1.05\text{ mg/L}.\end{align*}\textbf{Part (b): Critical Travel Time $t_c$:}\\From the derivative of the Streeter-Phelps oxygen deficit equation:\begin{equation*}t_c = \frac{1}{k_r - k_d} \ln\left\{ \frac{k_r}{k_d} \left[ 1 - D_0 \frac{k_r - k_d}{k_d L_0} \right] \right\}.\end{equation*}Substituting parameters ($k_r = 0.60, k_d = 0.25, k_r - k_d = 0.35$):\begin{align*}\frac{k_r}{k_d} &= \frac{0.60}{0.25} = 2.40 \\\frac{D_0(k_r - k_d)}{k_d L_0} &= \frac{1.05 \times 0.35}{0.25 \times 20.70} = \frac{0.3675}{5.175} = 0.071014 \\1 - 0.071014 &= 0.928986 \\\text{Term inside logarithm} &= 2.40 \times 0.928986 = 2.22957 \\\ln(2.22957) &= 0.80181 \\t_c &= \frac{0.80181}{0.35} = 2.2909\text{ days} \approx 2.29\text{ days}.\end{align*}\textbf{Part (c): Critical Deficit $D_c$ and Minimum DO:}\\At the critical point, rate of deoxygenation equals rate of reaeration ($k_d L(t_c) = k_r D_c$):\begin{equation*}D_c = \frac{k_d}{k_r} L_0 e^{-k_d t_c} = \frac{0.25}{0.60} \times 20.70 \times e^{-0.25 \times 2.2909} = 8.625 \times e^{-0.5727} = 8.625 \times 0.5640 = 4.864\text{ mg/L}.\end{equation*}The minimum dissolved oxygen is:\begin{equation*}\text{DO}_{\text{min}} = \text{DO}_{\text{sat}} - D_c = 8.80 - 4.864 = 3.936\text{ mg/L} \approx 3.94\text{ mg/L}.\end{equation*}\textbf{Regulatory Evaluation:} The CPCB Class D aquatic wildlife criterion requires $\text{DO} \ge 4.0\text{ mg/L}$. Because $3.94\text{ mg/L} < 4.0\text{ mg/L}$, the river experiences severe hypoxia, causing fish kills and endangering benthos. Wastewater treatment to reduce influent BOD is legally mandatory.\\\textbf{Part (d): Downstream Distance:}\\\begin{equation*}x_c = u \times t_c = 64.8\text{ km/day} \times 2.2909\text{ days} = 148.45\text{ km}.\end{equation*}The critical sag nadir is reached $148.45\text{ km}$ downstream of the discharge outfall.
\end{solbox}

\begin{examinertip}
\textbf{Marking Scheme Breakdown (Total: 10 Marks):}\\
\textbullet\ \textbf{3 Marks:} Clear formulation of governing environmental identities, equations, and thermodynamic/constitutional foundations.\\
\textbullet\ \textbf{4 Marks:} Rigorous algebraic/numerical substitution, dimensional consistency, and unit handling.\\
\textbullet\ \textbf{3 Marks:} Critical evaluation of environmental implications, regulatory criteria (CPCB/UNDRR/IUCN), and policy mitigations.
\end{examinertip}
